\chapter*{Abstract}
\addcontentsline{toc}{chapter}{Abstract}

Chemistry education contains foundational concepts that are often perceived as difficult for students to understand, particularly in terms of their application to real-world situations. This difficulty may lead to cognitive overload, especially when instruction lacks opportunities for practical engagement. Interactive learning approaches allow students to actively apply and develop their analytical skills while discovering how chemistry concepts relate to everyday experiences.

This study explores the use of Game-Based Learning (GBL) through the development of [GAME NAME HERE], an educational serious game designed as a learning material for Grade 6 students under the MATATAG curriculum in the Philippines. The game focuses on physical and chemical changes of matter, expanding students' understanding beyond the three common states of matter. Through interactive gameplay, students engage with concepts such as phase changes and heat transfer while learning to distinguish physical changes from chemical changes.

The primary goal of this study is to improve students' conceptual understanding, motivation, and emotional engagement by reinforcing hands-on learning and real-world applications of chemistry concepts. This project aims to provide both students and teachers with an accessible and engaging instructional resource that supports visualization, experimentation, and the development of analytical skills in basic chemistry education.

\vspace{1cm}

\noindent\textbf{Keywords:} physical and chemical changes, game-based learning, educational serious game

\noindent\textbf{Sustainable Development Goals:} 4: Quality Education, 11: Sustainable Cities and Communities, 12: Responsible Consumption and Production
